\chapter{Swelling}

\section*{Definition}
Swelling (edema) is an abnormal buildup of fluid in the tissues of your body. It can affect any part of the body but is most noticeable in the hands, arms, feet, ankles, and legs. Swelling may be pitting (leaves an indent when pressed) or non-pitting (firm and does not indent), which helps your doctor identify the cause.

\section*{When to See a Doctor}
See your doctor if:
\begin{itemize}
  \item Swelling comes on suddenly in one leg, especially with pain, redness, or warmth (possible DVT)
  \item You have shortness of breath, chest pain, or trouble breathing (possible heart failure or pulmonary embolism)
  \item The swelling is in your face or throat and came on quickly (possible allergic reaction)
  \item You are pregnant and have sudden severe swelling in your hands or face
  \item The area around the swelling is red, warm, and tender (possible infection)
  \item You have decreased urination along with swelling (possible kidney problem)
\end{itemize}

\section*{How It Develops}
Swelling occurs when fluid leaks out of small blood vessels into the surrounding tissues faster than it can be drained away. This can happen when pressure inside blood vessels is too high, when proteins that hold fluid in the bloodstream are too low, or when the lymphatic system (which normally drains excess fluid) is blocked or damaged. Inflammation can also increase the permeability of blood vessels, allowing more fluid to escape.

\section*{Causes}
\begin{itemize}
  \item \textbf{Heart failure:} The heart cannot pump blood efficiently, causing fluid to back up in the legs and lungs
  \item \textbf{Kidney disease:} Damaged kidneys cannot remove excess fluid and salt, leading to swelling in the face, hands, and feet
  \item \textbf{Liver disease (cirrhosis):} Scarred liver produces less protein, and fluid builds up in the abdomen and legs
  \item \textbf{Venous insufficiency:} Weak valves in leg veins allow blood to pool, causing ankle and leg swelling
  \item \textbf{Deep vein thrombosis (DVT):} A blood clot in a deep leg vein causes sudden swelling, pain, and warmth in one leg
  \item \textbf{Lymphedema:} Blocked or damaged lymphatic vessels prevent fluid drainage, often after surgery or radiation
  \item \textbf{Medication side effects:} Calcium channel blockers, NSAIDs, steroids, and some diabetes medications can cause swelling
  \item \textbf{Pregnancy:} Mild swelling is common, but sudden severe swelling may signal preeclampsia
  \item \textbf{Infection or inflammation:} Cellulitis, arthritis, or injury can cause localized swelling
  \item \textbf{Allergic reaction:} Swelling of the face, lips, or throat can be a medical emergency
\end{itemize}

\section*{Related Symptoms}
\begin{itemize}
  \item Puffiness or tightness of the skin
  \item Skin that holds an indent when pressed (pitting edema)
  \item Shiny or stretched skin over the swollen area
  \item Weight gain from retained fluid
  \item Shortness of breath (if fluid builds up in the lungs)
  \item Reduced flexibility in nearby joints
  \item Pain or heaviness in the affected limb
\end{itemize}
\textbf{Important:} Sudden swelling in one leg with pain and warmth may be a blood clot and requires urgent medical attention.

\section*{Medical Evaluation}
\begin{itemize}
  \item Your doctor will ask when the swelling started, whether it is in one area or throughout your body, and what makes it better or worse.
  \item They will check for pitting by pressing on the swollen area and measure the extent of the swelling.
  \item They will listen to your heart and lungs and check your abdomen for fluid buildup.
  \item Blood tests may include kidney and liver function, protein levels, and a blood count.
  \item Imaging such as ultrasound can check for blood clots or evaluate blood flow in the affected area.
  \item Urine tests can check for protein loss that might indicate kidney disease.
\end{itemize}

\section*{Treatment Options}
\begin{itemize}
  \item Treatment is directed at the underlying cause — heart failure, kidney disease, and liver disease each require specific management.
  \item Diuretics are used only for selected systemic causes and require monitoring; they are not a general treatment for unexplained swelling.
  \item Compression is considered for selected venous or lymphatic swelling only after arterial circulation and the cause have been assessed.
  \item Blood thinners are used to treat blood clots (DVT).
  \item Elevation and movement may help selected dependent or venous swelling; dietary changes should be individualized.
  \item Manual lymphatic drainage (a special type of massage) can help with lymphedema.
\end{itemize}

\section*{Self-Care and Home Management}
\begin{itemize}
  \item Elevate the swollen area above heart level when resting to help fluid drain.
  \item Follow individualized dietary advice and do not impose strict salt or fluid restriction without medical guidance.
  \item Wear compression stockings if recommended by your doctor for leg swelling.
  \item Move around regularly — sitting or standing still for long periods can worsen swelling.
  \item Drink a normal amount of fluid unless a clinician has prescribed a different plan.
  \item Avoid tight clothing or jewelry that can restrict circulation around the swollen area.
\end{itemize}

\section*{References}
\begin{thebibliography}{99}
\bibitem{swelling:ref1} Scallan J, Huxley VH, Korthuis RJ. ``Capillary fluid exchange: Regulation, functions, and pathology.'' \textit{Synthesis Lectures on Integrated Systems Physiology}, 2010;2(1):1--94.
\bibitem{swelling:ref2} Wakefield TW, Greenfield LJ. ``Diagnostic approaches to deep vein thrombosis.'' \textit{CHEST}, 2000;118(1):225--232.
\bibitem{swelling:ref3} Guyton AC, Hall JE. ``Local control of blood flow by the tissues, and humoral regulation.'' In: \textit{Textbook of Medical Physiology}, 14th ed. Elsevier, 2020.
\bibitem{swelling:ref4} Simini B. ``Acute limb swelling: A practical approach to diagnosis and management.'' \textit{Journal of the Royal College of Physicians of Edinburgh}, 2021;51(2):147--153.
\bibitem{swelling:ref5} Kearon C, Akl EA, Ornelas J, et al. ``Antithrombotic therapy for VTE disease: CHEST guideline and expert panel report.'' \textit{CHEST}, 2016;149(2):315--352.
\bibitem{swelling:ref6} Trayes KP, Studdiford JS, Pickle S, et al. ``Edema: Diagnosis and management.'' \textit{American Family Physician}, 2013;88(2):102--110.
\end{thebibliography}
